This document is the theoretical foundation of the RecoCrysp tutorial series.
It develops, from first principles, the physics and mathematics of Positron
Emission Tomography (PET) image reconstruction as implemented in the
\code{RecoCrysp} Julia library: the scanner geometry and line-of-response
model, the voxel grid, the Joseph forward/back projectors and their matched
adjoint, the detector point-spread function, the complete physical data model
(attenuation, normalization, scatter and randoms), the Poisson statistical
model, and maximum-likelihood iterative reconstruction (MLEM and OSEM),
including noise behaviour and semi-convergence. Each subsequent tutorial
(\code{tutorial\_example1}, \dots) applies this material to a concrete, fully
worked example. The treatment is non-time-of-flight and listmode-oriented, and
runs identically on CPU and GPU (including Apple Metal).
